\section{其他养殖与水产品：肉禽、蛋禽、牛羊、水产}
\label{sec:other-livestock}

生猪之外的养殖品种（肉牛、肉羊、白羽肉鸡、蛋鸡、水产）没有期货合约，
但它们构成了中国农业产值的另一大半，而且\hl{与猪周期共享同一条成本曲线、
同一种产能机制、同一批消费者}。本章的问题是：这些品种的周期是猪周期的“变奏”，
还是各自独立的周期？答案取决于一个变量——\emph{产能调整所需的时间}。

\subsection{产能时间轴：不同养殖品种的周期长度}
\label{sec:other-biology}

\begin{center}
\begin{tikzpicture}[font=\scriptsize, x=0.0082cm, y=1cm]
  % 时间轴：0 ~ 52 个月（200 刻度单位 = 6 个月）
  \draw[-{Stealth[length=4pt]}, InkGray, line width=0.8pt] (0,0) -- (1760,0);
  \foreach \m/\lb in {0/0, 200/6, 400/12, 600/18, 800/24, 1100/33, 1400/42}
    \draw[InkGray] (\m,0) -- (\m,-0.12) node[below, font=\tiny] {\lb};
  \node[font=\tiny, InkGray, anchor=west] at (1620,-0.30) {个月};
  % 各类品种的产能响应条
  \draw[AgriGold, line width=7pt, line cap=round] (0,0.70) -- (400,0.70);
  \node[font=\scriptsize, anchor=west, AgriGold!70!black] at (0,1.02) {蛋鸡、肉鸡：约 6---12 个月};
  \draw[HogPink, line width=7pt, line cap=round] (0,1.55) -- (550,1.55);
  \node[font=\scriptsize, anchor=west, HogPink!70!black] at (0,1.87) {生猪：约 10---12 个月};
  \draw[AquaTeal, line width=7pt, line cap=round] (0,2.40) -- (900,2.40);
  \node[font=\scriptsize, anchor=west, AquaTeal!70!black] at (0,2.72) {肉羊：约 18---24 个月};
  \draw[OilBlue, line width=7pt, line cap=round] (0,3.25) -- (1650,3.25);
  \node[font=\scriptsize, anchor=west, OilBlue!80!black] at (0,3.60) {肉牛：约 24---44 个月以上};
  % 参照线：生猪的产能周期终点
  \draw[InkGray, dashed, line width=0.5pt] (550,-0.1) -- (550,2.9);
\end{tikzpicture}
\end{center}

这张图解释了本章最核心的判断：\hl{产能时间越长的品种，价格周期越长、振幅越大、
与猪周期的错位越明显}。肉牛从犊牛到出栏约 24 个月，从补栏新生母牛到育成牛出栏
通常需 44 个月以上，且牛为单胎动物\cite{ndrc2026beef}——
这意味着 2023 年开始的母牛去化要到 2027 年之后才能反映为供给收缩；
而蛋鸡只需淘汰与补栏，去产能可在半年内完成。
因此当四个品种同时处于下行期时，它们的\emph{底部时间完全不同}。

\subsection{肉牛：政策触发的产能出清}
\label{sec:other-beef}

肉牛是本章中周期最完整、也最典型的外生冲击案例。价格自 2022 年末持续下降至 2025 年一季度，
2025 年二季度起止跌转涨，结束了持续两年的下行周期\cite{ndrc2026beef}。
产能方面，肉牛存栏于 2023 年三季度达到 \hl{10524 万头的峰值}，
随后持续下降，2025 年一季度降至 9462 万头（2020 年二季度以来最低），
2025 年末全国牛存栏 9608 万头（同比 -4.37\%），
2026 年二季度末进一步下滑至 9557 万头（同比 -4.3\%）\cite{ndrc2026beef,caas2026agri}。
同期出栏量与产量仍在增长（2025 年出栏 5133 万头、牛肉产量 801.00 万吨、创近五年新高），
说明\hl{当前的产量增长来自“产能去化后的集中出栏”}，是不可持续的。

价格已经开始反映供给收缩：2025 年肉牛交易均价 26.45 元/公斤（同比 +16.68\%），
国产鲜牛肉批发均价 59.53 元/公斤（+17.56\%）\cite{caas2026agri}；
2026 年 8 月末至 9 月初牛肉批发市场周均价 69.59 元/公斤（同比高 6.8\%）\cite{xinhua2026week}。
更关键的政策变量出现在 2025 年 12 月 31 日：商务部公告裁定进口牛肉数量增加
对中国国内产业造成严重损害，决定以“国别配额及配额外加征关税”形式实施保障措施，
期限 3 年（2026 年 1 月 1 日---2028 年 12 月 31 日），配额外加征关税税率 \hl{55\%}
\cite{mofcom2026beef}。这将进口牛肉（此前八件套均价约 48.64 元/公斤）
的价格优势大幅削弱，\emph{等于给国内肉牛价格加上了一个政策底}。

\subsection{肉羊：被猪价拖累的独立周期}
\label{sec:other-mutton}

肉羊提供了本章中最有趣的“交叉抑制”样本。产能上，经过两年多的调整，
2025 年第一季度末全国肉羊存栏 3.00 亿只，处于 2017 年以来本轮周期的低位；
2024 年新生羔羊数量同比下降 4.8\%，2025 年 1---5 月继续下降 3.36\%\cite{lj2026mutton}。
按产能逻辑，价格应当上涨，但 2025 年羊肉价格只走出“先抑后扬”的温和反弹：
1 月 70.1 元/kg 降至 7 月 68.9 元/kg，再回升至 12 月 71.6 元/kg，全年仅上涨 2.1\%
\cite{lj2026mutton}。

原因在于替代品的拖累：2025 年猪肉集贸市场均价同比下降 \hl{18.9\%}、
鸡肉价格下跌 5.8\%，作为羊肉核心替代品的牛肉价格虽上涨 8.5\% 但整体仍处近五年相对低位，
对羊肉价格形成明显的\hl{交叉抑制效应}\cite{lj2026mutton}。
到 2026 年，随着猪价企稳与牛羊供给收缩共振，羊肉批发价在 2026 年第 36 周
达到 66.13 元/公斤（同比高 8.9\%）\cite{xinhua2026week}。
\hl{肉羊的案例说明：产能周期决定方向，替代品价格决定幅度}。

\subsection{白羽肉鸡与蛋鸡：短周期、快反转}
\label{sec:other-poultry}

禽类的产能调整最快，因此周期最短、反转最急。白羽肉鸡在 2025 年经历了
“祖代更新创历史新高 + 出栏价格创近十年低位 + 产业链利润断崖”的组合：
祖代年度更新 157.41 万套（+4.89\%，超过 2013 年历史峰值），
其中海外引种占 39.92\%、国内自繁品种占 60.08\%；
在产父母代种鸡 5604.58 万套（+11.86\%），连续 3 年创历史新高；
全年出栏均价仅 7.28 元/kg（同比 -6.60\%，近 10 年低位）\cite{yangji2026chicken}。
出栏量却仍在增长：2025 年白羽肉鸡出栏 90.57 亿只（+5.97\%），
鸡肉产量 1930.19 万吨（+8.50\%），全国鸡肉产量 2803.44 万吨（+6.18\%）\cite{yangji2026chicken}。

反转发生在 2026 年：1---4 月毛鸡平均价格 3.63 元/斤（同比 +4.31\%），
毛鸡养殖利润 0.71 元/只（同比 +\hl{297.22\%}），
祖代引种因海外引种中断而全部转为国内自繁（更新量 29.38 万套、同比 -13.84\%）
\cite{yangji2026chicken}。\emph{引种中断是一个被低估的供给冲击}——
祖代供给的下滑将在 14 个月后传导到父母代、再传导到商品代。

蛋鸡的周期与生猪高度同步但幅度更大：2025 年全国在产蛋鸡年均存栏约 12.2 亿羽（+1.3\%），
9 月峰值突破 12.7 亿羽，全年平均淘汰日龄预计超过 520 天（去产能缓慢），
鸡蛋每千克平均亏损约 0.93 元\cite{aige2025layer}。
2026 年在产蛋鸡存栏逐月下降，8 月约 12.82 亿羽（同比 -6.08\%），
主产区蛋价于 2026 年 5 月底冲至 5 元/斤、创近十年同期新高\cite{zhuochuang2026layer}；
2026 年第 36 周鸡蛋批发价 10.83 元/公斤（同比 +\hl{41.9\%}）\cite{xinhua2026week}，
上半年禽蛋产量 1690 万吨（同比 -2.2\%）\cite{stats2026h1}。
在第 \ref{sec:corr-transmission} 节的传导检验中，\emph{鸡蛋是唯一稳健的、
与生猪周期同源的品种}，这与其“同样的产能机制 + 更短的调整时间”完全一致。

\subsection{水产：产能周期之外的另一个周期}
\label{sec:other-aqua}

水产品与前四类有本质区别：\hl{它的供给由养殖面积与投苗量决定，而不是由存栏决定}，
且需求端由居民蛋白消费结构升级驱动，与猪价的关系更弱
（第 \ref{sec:corr-sw} 节中水产养殖指数与猪价呈显著反周期）。

规模上，2025 年全国水产品总产量 7652.14 万吨（同比 +4.00\%），
其中养殖产量 6325.62 万吨（+4.38\%）、捕捞产量 1326.52 万吨（+2.23\%），
养殖与捕捞之比为 82.7:17.3\cite{moa2026fishery}；
水产品进口量 707.47 万吨（+2.05\%），贸易逆差扩大至 42.18 亿美元\cite{moa2026fishery}。

价格周期由 2024 年的\hl{投苗量大幅下降}启动：2025 年二季度起淡水鱼价格大幅上涨，
草鱼塘口均价 8.5 元/斤（同比 +\hl{43\%}）、价格高点达 14 元/公斤并创历史同期新高
\cite{nfnews2025grasscarp}；而海水产品同期走弱（南美白对虾 -8.98\%）\cite{caas2026agri}。
但到 2026 年第 36 周，淡水鱼价格全面回落：鲤鱼 14.12、草鱼 17.85、鲫鱼 20.22 元/公斤，
环比均小幅下跌\cite{xinhua2026week}。
\hl{这构成一个典型的“投苗$\to$上市”滞后周期（约 12---18 个月），
与生猪的“配种$\to$出栏”周期同构，但时间更长、弹性更大}。

\subsection{养殖业的周期矩阵}
\label{sec:other-matrix}

\begin{table}[htbp]
\centering
\small
\begin{tabular}{@{}lrrrrl@{}}\toprule
品种 & 产能周期 & 2025 年状态 & 2026 年方向 & 与猪周期 & 关键变量 \\
\midrule
生猪   & 10$\sim$12 月 & 深度亏损、政策去化 & 磨底回升 & 基准      & 能繁母猪、政策 \\
肉牛   & 24$\sim$44 月 & 存栏 -4.4\%、价格上涨 & 上行      & 错位      & 进口保障措施 \\
肉羊   & 18$\sim$24 月 & 存栏低位、价格微涨 & 温和上行  & 被抑制    & 替代品价格 \\
白羽肉鸡 & 6$\sim$12 月 & 价格近十年低位、利润崩塌 & 快速反转 & 高度同步  & 祖代引种 \\
蛋鸡   & 6$\sim$10 月 & 存栏高位、全面亏损 & 快速反转  & 高度同步  & 淘汰日龄 \\
水产   & 12$\sim$18 月 & 投苗下降、价格冲高 & 高位回落  & 反周期    & 投苗量、进口 \\
\bottomrule
\end{tabular}
\caption{养殖业各品种的周期特征对照（2025---2026）}
\label{tab:other-livestock-matrix}
\srcfile{数据来自本章引用的农业农村部、国家统计局、商务部公告与各产业技术体系报告。}
\end{table}

\begin{keybox}[本章三个结论]
\normalsize
\textbf{① 产能时间决定周期形态。}肉牛的 44 个月产能刚性使其供给缺口无法在短期内回补，
价格上行一旦形成就具有持续性；禽类的半年产能使其价格反复快速反转，
“高价$\to$亏损$\to$去化$\to$高价”的循环在两年内可以完成两轮。\par\vspace{3pt}
\textbf{② 替代效应是养殖业内部的传导机制。}2025 年猪价下跌 18.9\% 直接压制了羊肉价格，
使肉羊在自身产能已去化的条件下仅录得 2.1\% 的涨幅。\hl{养殖品种之间的相关性通过
“消费替代”而非“成本传导”实现}，这与第 \ref{sec:corr} 章商品价格相关性的结论形成互补。\par\vspace{3pt}
\textbf{③ 政策在养殖业中的作用从“托底”走向“出清”。}生猪产能调控方案、
进口牛肉保障措施、祖代引种中断，都是行政力量直接改变供给曲线的例子。
在 2025---2026 年的农业周期中，\hl{政策是产能去化的加速器，而不是价格的保证}。
\end{keybox}
